\chapter{A Concrete Naming Certificate for $\CompactIntervalInverseModulusUp$}
\label{ch:concrete-instances-compact-interval-inverse-modulus-namecert}
\origin{human}

Following Bishop--Bridges constructive analysis on inverse functions over compact real intervals, the $\CompactIntervalInverseModulusUp$ packet records the finite modulus data needed to read a monotone inverse branch without invoking host inverse functions or quotient real equality. It sits over the compact interval route of \autoref{ch:concrete-instances-closed-bounded-interval-compact-namecert}, the monotone inverse branch of \autoref{ch:concrete-instances-monotone-inverse-namecert}, the uniform-modulus consumer of \autoref{ch:concrete-instances-uniformmodulus-namecert}, the finite window surface of \autoref{ch:concrete-instances-streamname-namecert}, the regular rational readback of \autoref{ch:concrete-instances-regseqrat-namecert}, the dyadic tolerance ledger of \autoref{ch:concrete-instances-dyadicratcore-namecert}, and the terminal real-name seal of \autoref{ch:concrete-instances-real-namecert}. The packet is a compact-interval inverse-modulus certificate for Real and Completion consumers; it is not an open-cover compactness theorem, a trichotomy principle, a selected inverse function, or an ambient $\RealUp$ completeness assertion.

\begin{definition}[Compact interval inverse modulus carrier]
\label{def:compact-interval-inverse-modulus-carrier}
A $\CompactIntervalInverseModulusUp$ carrier is a finite $\BHist$ packet
$$
I=(K,M,F,L,W,D,R,S,H,C,P,N)
$$
with the following displayed rows:
\begin{enumerate}
\item $K$ is the $\ClosedBoundedIntervalCompactUp$ compact interval source row, exposing the closed bounded interval, finite-net compactness route, compact-metric row, and endpoint $\RealUp$ seal.
\item $M$ is the $\MonotoneInverseUp$ branch row, naming the monotone inverse branch only through displayed source, target, and order rows.
\item $F$ is the finite-net row over $K$, listing the net centers, interval cells, image comparisons, and radius data inspected by the inverse-modulus route.
\item $L$ is the inverse-modulus ledger, recording how target separation requests are discharged by finite source cells, monotone branch data, and compact interval bounds.
\item $W$ stores the finite $\StreamNameUp$ windows opened on source endpoints, net centers, branch values, and target requests.
\item $D$ stores the $\DyadicRatCoreUp$ tolerance rows for endpoint radii, image gaps, requested target tolerances, and residual readback error.
\item $R$ is the $\RegSeqRatUp$ readback row produced from $W$ and $D$ before any real-facing seal is exposed.
\item $S$ is the terminal $\RealUp$ seal row for the source interval and the inverse-modulus readback.
\item $H$ records componentwise $\hsame$ transport for the compact interval source, monotone inverse branch, finite-net row, inverse-modulus ledger, stream windows, dyadic tolerance rows, regular readback, Real seal, provenance, and local naming row.
\item $C$ records displayed $\Cont$ replay from compact interval inspection through monotone branch routing, finite-net handoff, tolerance accounting, regular readback, and Real sealing.
\item $P$ is the $\Pkg$ provenance over $\CompactIntervalInverseModulusUp$, $\ClosedBoundedIntervalCompactUp$, $\MonotoneInverseUp$, $\UniformModulusUp$, $\StreamNameUp$, $\DyadicRatCoreUp$, $\RegSeqRatUp$, $\RealUp$, $\BHist$, $\hsame$, $\Cont$, and $\NameCert$.
\item $N$ is the local $\NameCert_{\CompactIntervalInverseModulusUp}$ row.
\end{enumerate}
The classifier compares accepted packets only through $K,M,F,L,W,D,R,S,H,C,P,N$. It has no coordinate for a host inverse function, quotient real equality, trichotomy, countable choice, open-cover compactness, or ambient $\RealUp$ completeness.
\end{definition}

\begin{theorem}[Compact interval inverse modulus NameCert obligations]
\label{thm:compact-interval-inverse-modulus-namecert-obligations}
For every accepted $\CompactIntervalInverseModulusUp$ carrier $I=(K,M,F,L,W,D,R,S,H,C,P,N)$, the local $\NameCert_{\CompactIntervalInverseModulusUp}$ surface consists exactly of carrier admission, compact interval exposure, monotone inverse branch exposure, finite-net handoff, inverse-modulus ledger exactness, $\StreamNameUp$ window visibility, $\DyadicRatCoreUp$ tolerance accounting, $\RegSeqRatUp$ readback, $\RealUp$ sealing, componentwise $\hsame$ transport, displayed $\Cont$ replay, $\Pkg$ provenance, and local naming.
\end{theorem}
\begin{proof}
Project the carrier of \autoref{def:compact-interval-inverse-modulus-carrier}. Its analytic rows are $K,M,F,L,W,D,R,S$, and its structural rows are $H,C,P,N$.

The compact interval source is upstream. A consumer must first read $K$, then the monotone branch row $M$, before the finite-net row $F$ and inverse-modulus ledger $L$ can be inspected.

The finite real-name rows are downstream of that ledger. The stream windows $W$ and dyadic tolerance rows $D$ feed the regular readback $R$, and the real-facing seal $S$ is reached only after that finite readback is displayed.

Finally $H,C,P,N$ preserve, replay, source, and name the same packet. Since the refused host objects are not carrier coordinates, none can appear as a local NameCert obligation.
\end{proof}

\begin{theorem}[Compact interval inverse modulus finite-net handoff]
\label{thm:compact-interval-inverse-modulus-finite-net-handoff}
Every $\UniformModulusUp$ consumer of an accepted $\CompactIntervalInverseModulusUp$ packet factors through the displayed finite route
$$
\ClosedBoundedIntervalCompactUp
\longmapsto
\MonotoneInverseUp
\longmapsto
F,L
\longmapsto
W,D,R,S
\longmapsto
H,C,P,N .
$$
The handoff supplies compact interval data, monotone inverse branch data, finite-net cells, inverse-modulus ledger rows, stream windows, dyadic tolerances, regular readback, and Real sealing as one finite packet. It exports no host inverse function, quotient real equality, trichotomy, countable choice, open-cover compactness, or ambient Real completeness.
\end{theorem}
\begin{proof}
Start with the compact interval source and branch rows. The rows $K$ and $M$ force any inverse-modulus read to expose the compact interval certificate and the monotone inverse branch before a modulus request is evaluated.

Next read the finite-net and ledger rows. The row $F$ lists the finite cells and image comparisons, while $L$ records the inverse-modulus accounting used by the downstream $\UniformModulusUp$ consumer.

Then attach the finite real-name handoff. The stream windows $W$ and dyadic tolerances $D$ produce the regular readback $R$, and the terminal seal $S$ is available only after that readback.

The structural rows $H,C,P,N$ preserve this order. Hence the consumer receives a finite inverse-modulus route, not a host inverse map or an ambient compactness principle.
\end{proof}

\closureat{CompactIntervalInverseModulusUp}{seedStr}

\begin{closurestatus}{\CompactIntervalInverseModulusUp}
  \constructivestory{A $\CompactIntervalInverseModulusUp$ packet is generated from a compact interval source, monotone inverse branch, finite-net row, inverse-modulus ledger, StreamName windows, DyadicRatCore tolerance rows, RegSeqRat readback, RealUp seal, componentwise hsame transport, displayed Cont replay, Pkg provenance, and a local NameCert row.}
  \theoryclosure{\seedClosure}
  \scopeclosed{\autoref{def:compact-interval-inverse-modulus-carrier}, \autoref{thm:compact-interval-inverse-modulus-namecert-obligations}, and \autoref{thm:compact-interval-inverse-modulus-finite-net-handoff} seed the compact-interval inverse-modulus carrier and its finite-net handoff surface.}
  \formalstatus{\unformalizedV}
  \bridgestatus{none}
  \notclaimed{This chapter does not close host inverse functions, quotient real equality, trichotomy, countable choice, open-cover compactness, ambient RealUp completeness, Lean verification, or bridge checking.}
  \upgradepath{Obligation closure requires checked carrier admission, compact interval exposure, monotone branch exposure, finite-net handoff, inverse-modulus ledger exactness, finite real-name readback, transport, replay, provenance, and local naming for BEDC.Derived.CompactIntervalInverseModulusUp.}
\end{closurestatus}
